\chapter{Introduzione}
\label{cap:td-introduzione}

\section{Scopo del documento}

Questo documento descrive come Tech4All è verificato: la strategia adottata,
il metodo con cui i casi di test sono stati derivati, la loro specifica, i
risultati dell'esecuzione e i limiti della verifica.

Rispetto alla struttura suggerita da IEEE 829~\cite{ieee829}, il documento
riunisce in un solo testo ciò che lo standard distribuisce fra piano di test,
specifica dei casi, rapporto sugli incidenti e rapporto riassuntivo. La
ragione è pratica: quei quattro documenti condividono premesse, glossario e
strumenti, e mantenerli separati per un sistema di questa scala produce
ripetizione senza aggiungere rigore.

La riunione dei quattro documenti ha però una conseguenza sull'ordine di
lettura: il documento è redatto man mano che la verifica procede, e i
capitoli che descrivono strategia e casi di test precedono l'esecuzione
anche cronologicamente. Il \cref{cap:td-esecuzione} riporta l'esito
dell'ultima esecuzione della suite e i difetti individuati fino a quel
momento.

\section{Relazione con gli altri documenti}

Il documento chiude la catena di tracciabilità del progetto.

\begin{description}[leftmargin=1.4cm,style=nextline]
  \item[Dal RAD] provengono i requisiti funzionali e non funzionali, che sono
    l'oggetto della verifica, e i criteri di accettazione, che ne sono
    l'oracolo. Un requisito non funzionale del RAD è formulato in modo da
    poter diventare un'asserzione: è la ragione per cui, in quel documento,
    ogni requisito ha un criterio osservabile.
  \item[Dal SDD] provengono i confini delle unità da verificare: servizi,
    middleware e serializzazione. La scelta di iniettare le dipendenze dai
    costruttori, motivata in quel documento, è ciò che rende la verifica in
    isolamento possibile senza artifici.
  \item[Verso entrambi] questo documento riporta, in
    \cref{sec:td-limiti}, ciò che la verifica \emph{non} copre: sono i punti
    in cui i compromessi dichiarati nel SDD si traducono in rischio residuo.
\end{description}

La matrice che lega ogni requisito ai casi di test che lo verificano è
nell'appendice~\ref{cap:td-app-tracciabilita}.

\section{Riferimenti}

Il metodo di derivazione dei casi di test è quello della \term{category
partition} di Ostrand e Balcer~\cite{ostrand1988}. La terminologia relativa
alla documentazione della verifica segue IEEE 829~\cite{ieee829}. Le
classi di vulnerabilità citate nei casi di test di sicurezza fanno
riferimento alla classificazione OWASP~\cite{owasp2021}.
